The two figures below are the descriptive pass that precedes the Bayesian analysis, run on pure sinusoids swept from $2$ to $250$\,Hz in one-hertz steps at $f_s=512$\,Hz over the five geometries of the initial design. They are reported here rather than in the body because neither is a hypothesis test: both average over phase, use a single signal family, and carry no statement of uncertainty. Their role is to show that the phenomenon the models of \nameref{sec:four} are built to measure is present in the raw sweep. The summary statistics read off them are given in \nameref{sec:seven}.

\begin{figure*}[!p]
    \centering
    \includegraphics[width=0.92\textwidth]{figures/FIG_ALL_reconstruction.png}
    \caption{Behavioural reading. Each panel plots, for one geometry, the dominant frequency of the recursive multi-step forecast against the frequency of the tone supplied as context, averaged over the non-redundant starting phases with one standard deviation shaded. The tone is given as context, the model's median forecast is appended to that context and the model re-invoked, and the loop continues until a full window has been generated; the dominant frequency is then read as the largest magnitude of the discrete Fourier transform after the mean is removed. The recursion is imposed by this procedure and is not the model's own generation mode, since Chronos-Bolt emits its whole horizon in one step. A point on the dotted diagonal is a tone rebuilt correctly; a flat segment is a set of distinct inputs collapsing onto one output; a jump is a tone the model rebuilds that was never supplied. Dashed verticals are the predicted candidates $\mathcal F_{\mathrm{lock}}$ of that geometry. The reading is behavioural and records what the model outputs, not what its representation retains.}
    \label{fig:allReconstruction}
\end{figure*}

\subsection{Planned extension: scope}

Two additions are planned to this appendix. The first extends the behavioural reading of
\autoref{fig:allReconstruction} from the five geometries of the initial design to the fifteen of
\autoref{tab:hfModels}, so that the descriptive pass covers the same grid the inferential stage is
fitted on. The second is a direct spectral comparison between the signal supplied as context and
the signal the model returns, carried out on the three signal families of \nameref{sec:five}: a
pure tone, a Light TSMixup background and a KernelSynth background.

The comparison is run on the geometry $P=16$, $S=16$. That is the geometry of the published
Chronos-Bolt tokeniser, so the reading concerns the configuration a deployment would actually meet,
and since $P=S$ the two branches of $\mathcal F_{\mathrm{lock}}$ coincide in a single comb at
multiples of $f_s/16$, which leaves one predicted set to mark rather than two interlaced ones.

Its purpose is motivation rather than test. Light TSMixup and KernelSynth are the two synthetic
generators of the Chronos data-generation procedure, so a loss observed on them cannot be answered
with the objection that the probe is out of distribution; and the extent of that loss is what makes
the question the models of \nameref{sec:four} are built to decide worth asking in the first place.

\subsection{Planned extension: protocol}

For each signal family one realisation is produced at the probing length, the context of $480$
samples is supplied to the model, and the forecast of $64$ samples is retained at the median
quantile. Three spectra are then computed under the same window and the same mean removal: the
context, the true continuation over the $64$ samples that follow it, and the forecast over those
same $64$ samples. For the families built from sinusoidal sources the generating frequencies are
known by construction and are marked together with the candidate set $\mathcal F_{\mathrm{lock}}$ of
the geometry. The comparison of record is the forecast against the true continuation, since those
two share a window and therefore a frequency resolution; the context spectrum is reported beside
them as a statement of what entered the model, at its own finer resolution.

\subsection{Planned extension: what it can and cannot show}

Three limits fix what may be read off the comparison, and each is a property of the measurement
rather than of the model. A $64$-sample window at $f_s=512$\,Hz resolves $8$\,Hz, so the comb of the
reference geometry falls on bins but a candidate and its controls at $f_k\pm\delta$ do not separate:
the picture can show loss, not the paired contrast the models estimate. The forecast retained is a
median quantile, and a conditional median is smoother than a sample path, so part of any
high-frequency attenuation belongs to the point forecast rather than to the tokenisation; the
matched triplet used in the body cancels that term, a raw spectrum does not. And KernelSynth draws
are broadband and non-stationary, so their spectra carry no lines to read against: that family
speaks to loss in general, while any site-level reading rests on Light TSMixup, whose components are
known one by one.

The comparison is therefore descriptive. It is evidence of frequency information loss, not evidence
of structural aliasing, and it is placed here so that the body may point to it as motivation without
any hypothesis resting on it.

% ------------------------------------------------------------------------------------------
% TO DO, once the figures are produced: insert them above and add the closing paragraph that
% states what the comparison showed on each of the three families. Keep the claim at the level
% of information loss; nothing here may be read as evidence of structural aliasing.
% ------------------------------------------------------------------------------------------

\begin{figure*}[!p]
    \centering
    \includegraphics[width=0.98\textwidth]{figures/FIG_ALL_spectral_recovery.png}
    \caption{Representational reading. Each strip is one geometry, and each cell one swept frequency, coloured by the cross-validated accuracy of a linear probe trained on the internal state at the output quantile head. The task is local and of constant difficulty at every frequency, whether a tone at $f-1$\,Hz can be separated from one at $f+1$\,Hz, and it holds out phases rather than frequencies, so no cell is inflated by a probe memorising individual frequencies. Green is accuracy at or above $0.9$, amber between $0.6$ and $0.9$, red below $0.6$, which is close to chance on a two-class task; grey marks a frequency the task does not cover, which is a statement about the design and not about the model. Dashed verticals are the predicted candidates of that geometry. Read against \autoref{fig:allReconstruction}, a frequency that is off the diagonal but still green has lost the tone at the projection rather than at the tokenisation; a red cell is the stronger condition that structural aliasing predicts.}
    \label{fig:allSpectralRecovery}
\end{figure*}
